%%%%%%%%%%%%%%%%%%%%%%%%%%%%%%%%%%%%%%%%%%%%%%%%%%%%%%%%%%%%%%%%%%%%%%
% Lớp 11 | Quan hệ vuông góc trong không gian | Khoảng cách
% Dạng 1H8?5-5: Đường vuông góc chung của hai đường thẳng chéo nhau
% Số câu: 2  (sinh tự động bởi tach_cau.py; từ nay thầy sửa trực tiếp file này)
%%%%%%%%%%%%%%%%%%%%%%%%%%%%%%%%%%%%%%%%%%%%%%%%%%%%%%%%%%%%%%%%%%%%%%

%%%%%%%%%%%%%%%% TRẢ LỜI NGẮN %%%%%%%%%%%%%%%%

% nguồn: Ok_TraLoiNgan_Lop11_Lan1
\begin{ex}%[1H8V5-5]%[ma:1aa2b1be]
	Cho tứ diện đều ${A B C D}$ có cạnh ${4\sqrt{2}}$. Đoạn vuông góc chung  của hai đường thẳng chéo nhau ${A B, C D}$ bằng
	 \shortans{4}
	\loigiai{ 
		\begin{center}
			\begin{tikzpicture}[scale=0.8,font=\footnotesize,line join=round,line cap=round,>=stealth]
					 \path
						(0,0)  coordinate (B)
						++(-60:2)  coordinate (C)
						(6,0) coordinate (D) 
						 ($(D)!1/2!(C)$) coordinate (J)
						($(B)!2/3!(J)$) coordinate  (H)
						 ($(H)+(0,4)$) coordinate (A)
						 ($(A)!1/2!(B)$) coordinate (I);
					 \foreach \i in{B,C,D}{\draw (A)--(\i);};
					\draw (B)--(C)--(D) 
						 (A)--(J)
						(I)--(C); 
					\draw[dashed] (B)--(D) 
						(B)--(J) 
						 (D)--(I)--(J);
					 \pic[draw,angle eccentricity=1.8,angle radius=2mm]{right angle=B--J--C};
					\pic[draw,angle  eccentricity=1.8,angle radius=2mm]{right angle=B--I--C};
					\pic[draw,angle eccentricity=1.8,angle radius=2mm]{right angle=D--I--A}; 
					\pic[draw,angle eccentricity=1.8,angle radius=2mm]{right angle=A--J--D}; 
					 \foreach \i/\g in {A/90,B/180,C/-90,D/0,J/-90}
					 \fill[black] (\i) circle (1pt)+(\g:4mm) node[scale=1]{$\i$};
				\end{tikzpicture}
		\end{center}
		Gọi ${I, J}$  theo thứ tự là trung điểm của ${A B, C D}$.\\
		Các tam giác  ${A B C, A B D}$ đều có ${I}$ là trung điểm ${A B}$ nên
		 $\heva{&A B \perp C I \\
		&A  B \perp D I
		}  \Rightarrow A B \perp(I C D)$.
		Mà ${I J \subset(I C D) \Rightarrow A B \perp I J}$. \tagEX{1} 
		\noindent Tương tự, các tam giác ${A C D, B C D}$ đều có ${J}$ là trung điểm ${C  D}$ nên\\
		 $\heva{&
		CD\perp  AJ \\
		&CD\perp BJ 
		}\Rightarrow  CD\perp (ABJ)$.
		 Mà $J\subset (JAB)\Rightarrow CD\perp IJ$\tagEX{2}
		\noindent Từ (1) và  (2) suy ra ${I J}$ là đoạn vuông góc chung của hai đường thẳng ${A B, C D}$.\\
		Ta có ${C I=\dfrac{2 a \cdot \sqrt{3}}{2}=a \sqrt{3};  I J=\sqrt{C I^2-C J^2}=\sqrt{3 a^2-a^2}=a \sqrt{2}}$.\\
		Thay $a=\sqrt{8}$ thì d (${A B,  C D}$)= $4$.
	} 
\end{ex}

% nguồn: Ok_TraLoiNgan_Lop11_Lan2
\begin{ex}%[1H8V5-5]%[ma:a7cc5433]
	Cho hình chóp $S.A B C$, có $S A=S B=S C$, đáy là tam giác đều cạnh 1. Biết thể tích khối chóp $S.A B C$ bằng $\dfrac{\sqrt{3}}{3}$. Khoảng cách giữa hai đường thẳng $S A$ và $B C$ bằng (\textit{kết quả làm  tròn đến hàng phần trăm})
	 \shortans{0,86}
	 \loigiai{
		\begin{center}
			\begin{tikzpicture}[scale=0.8,font=\footnotesize,line join=round,line cap=round,>=stealth]
					\coordinate (A) at (-2,0); 
					 \coordinate (B) at (-1,-3);
					 \coordinate (C) at (3,0);
					\coordinate (M)  at ($(B)!0.5!(C)$);
					\coordinate (G) at ($(A)!2/3!(M)$); 
					 \coordinate (S) at ($(G)+(0,5)$);
					\coordinate (H)  at ($(S)!1/3!(A)$);
					 \draw (S)--(A)
						 (S)--(B)
						 (S)--(C)
						(B)--(C) 
						 (A)--(B);
					\draw[dashed,thin](A)--(C) 
						(S)--(G) 
						 (A)--(M)--(H);
					\foreach \i/\g in {S/90,A/180,B/-90,C/0,G/-90,M/-45,H/180}{\draw[fill=black](\i)  circle (1pt)
						($(\i)+(\g:3mm)$)  node[scale=1]{$\i$};}
				\end{tikzpicture}
		\end{center}
		Do hình chóp $S . A B C$ đều nên $S G$ là đường cao của hình chóp ($G$  là trọng tâm tam giác đều $A B C$).\\
		Kẻ $M H \perp S A$  tại $H$.\\
		 Ta có $\heva{&BC\perp AM\\&BC\perp SG}\Rightarrow BC\perp (SAM)$.\\
		Suy ra $BC\perp  MH$.\\
		Mà $MH\perp SA$ nên $M H$ là đoạn vuông góc chung của $S A$  và $B C$.\\
		Vậy khoảng cách giữa  hai đường thẳng $S A$ và $B C$ bằng $M H$.\\
		Ta có $V_{S.A B C}=\dfrac{1}{3}\cdot \dfrac{\sqrt{3}}{4} S G=\dfrac{\sqrt{3}}{3}  \Rightarrow S G=4$, $A G=\dfrac{\sqrt{3}}{3}$.\\
		$S A=\sqrt{A G^{2}+S  G^{2}}=\sqrt{\dfrac{3}{9}+16}=\dfrac{7 \sqrt{3}}{3}$.\\
		Ta có $S A\cdot M H=S G\cdot A M \Rightarrow M H=\dfrac{S  G\cdot A M}{S A}=\dfrac{3\cdot 4 \cdot \sqrt{3}}{2\cdot 7 \sqrt{3}}=\dfrac{6}{7} \approx 0{,}86$.
	 }
\end{ex}
